\chapter{Machine learning I: regressão, features e validação temporal}
\label{cap:ml1}

\abertura{Um modelo de aprendizado supervisionado aprende, com exemplos do passado, uma regra que liga o que se sabe hoje ao que acontece depois. Em mercados essa regra é fraca, muda com o tempo e se confunde com ruído. Por isso este capítulo usa dois modelos lineares simples e concentra o esforço no que decide se um resultado merece confiança: as variáveis, o alvo e a validação.}

\begin{objetivos}
  \item formular uma previsão financeira como problema supervisionado, datando cada variável;
  \item estimar e interpretar uma regressão linear e uma regressão logística;
  \item construir features de preço e volume usando só informação disponível na data;
  \item rotular observações por horizonte fixo e pelo método da tripla barreira;
  \item validar um modelo respeitando a ordem do tempo e reconhecer vazamentos;
  \item levar uma previsão até um backtest com custos e compará-la com referências simples.
\end{objetivos}

\section{Por que mercados são um problema diferente}

Reconhecer um gato em uma foto é um problema em que o aprendizado de máquina funciona muito bem: o rótulo é inequívoco, gatos não mudam de aparência de um ano para o outro e cada foto é um exemplo independente das demais. Prever o retorno de uma ação viola as três condições, e ainda acrescenta uma quarta dificuldade.

\begin{itemize}
  \item \textbf{Razão sinal-ruído muito baixa.} A parte previsível de um retorno é uma fração mínima da sua variação. \citet{gu2020} compararam mais de uma dezena de métodos na previsão do retorno mensal de ações americanas e os melhores, redes neurais, alcançaram um $R^2$ fora da amostra da ordem de 0,4\% ao mês. Mais de 99\% da variação ficou sem explicação, e ainda assim os autores mostram que esse pouco tem valor econômico.
  \item \textbf{Não estacionariedade.} A relação entre as variáveis muda com o regime de juros, com a liquidez e com a composição dos participantes. Um padrão estimado em uma década pode não existir na seguinte.
  \item \textbf{Poucos dados independentes.} Dezesseis anos de pregões são cerca de 4~mil observações diárias de um ativo, pouco para os padrões de aprendizado de máquina, e elas dependem umas das outras no tempo.
  \item \textbf{Adversários.} Um padrão lucrativo atrai capital, e as ordens de quem o explora movem os preços até que ele desapareça. Parte da baixa razão sinal-ruído vem dessa força de arbitragem.
\end{itemize}

Com tão pouco sinal, um bom ajuste ao passado se obtém quase sempre e informa pouco. O trabalho está em verificar se a regra sobrevive em dados que o modelo de fato não viu.

\subsection{O problema supervisionado e sua linha do tempo}

Em \termo{aprendizado supervisionado} cada exemplo é um par formado por um vetor de variáveis explicativas, as \termo{features}, e por uma resposta, o \termo{rótulo} ou alvo. O modelo é uma função $f$ ajustada para que $f(\vetor{x})$ fique perto de $y$ nos exemplos conhecidos. Se $y$ é um número, como um retorno, a tarefa é de \termo{regressão}; se é uma categoria, como ``sobe'' ou ``cai'', é de \termo{classificação}.

Com dados de mercado, o cuidado está nas datas. As features $\vetor{x}_t$ só podem conter o que se sabia no fechamento do dia $t$. O alvo é algo que acontece depois, por exemplo o log-retorno dos $h$ pregões seguintes,
\begin{equation}
  y_t = r_{t \to t+h} = \ln P_{t+h} - \ln P_t ,
  \label{eq:c09-alvo}
\end{equation}
e só fica conhecido em $t+h$. A Figura~\ref{fig:c09-linha} desenha duas observações consecutivas com $h=3$. Guarde dois fatos que voltam no capítulo inteiro: nenhuma seta de informação cruza a data $t$ da direita para a esquerda, e os rótulos de dias vizinhos se sobrepõem quando $h > 1$.

\begin{figure}[b]
\centering
\begin{tikzpicture}[x=8.2mm, y=1mm, font=\sffamily\footnotesize]
  \draw[-{Stealth[length=2mm]}] (0.2,0) -- (13.6,0);
  \foreach \i/\r in {1/t-7, 3/t-5, 5/t-3, 7/t-1, 8/t, 9/t+1, 10/t+2, 11/t+3, 12/t+4}
    \draw (\i,-1) -- (\i,1) node[above=0.5pt, font=\footnotesize] {$\r$};
  \foreach \i in {2,4,6} \draw (\i,-1) -- (\i,1);
  % observação t
  \fill[lilas] (1,-13) rectangle (8,-7);
  \draw (1,-13) rectangle (8,-7);
  \node at (4.5,-10) {features $\symbf{x}_t$: dados até o fechamento de $t$};
  \fill[roxo] (8,-13) rectangle (11,-7);
  \node[text=white] at (9.5,-10) {rótulo $y_t$};
  \node[anchor=east] at (0.9,-10) {obs.\ $t$};
  % observação t+1
  \fill[lilas] (2,-23) rectangle (9,-17);
  \draw (2,-23) rectangle (9,-17);
  \node at (5.5,-20) {features $\symbf{x}_{t+1}$};
  \fill[roxo] (9,-23) rectangle (12,-17);
  \node[text=white] at (10.5,-20) {rótulo $y_{t+1}$};
  \node[anchor=east] at (1.9,-20) {obs.\ $t+1$};
  % decisão e sobreposição
  \draw[densely dashed] (8,-1) -- (8,-30);
  \node[anchor=north, align=center] at (8,-30) {decisão\\de $t$};
  \draw[decorate, decoration={brace, amplitude=3pt, mirror}] (9,-25.5) -- (11,-25.5);
  \node[anchor=north, align=center] at (11.2,-27.5) {dois dias de retorno\\em comum};
\end{tikzpicture}
\caption{Linha do tempo de duas observações com horizonte $h=3$. As features da observação $t$ usam apenas o passado; o rótulo é o retorno dos três pregões seguintes e só existe em $t+3$. Os rótulos de $t$ e de $t+1$ compartilham dois dos três retornos diários.}
\label{fig:c09-linha}
\end{figure}

\section{Regressão linear}
\label{sec:c09-linear}

A regressão linear supõe que a resposta é uma soma ponderada das features mais um erro:
\begin{equation}
  y_i = \beta_0 + \beta_1 x_{i1} + \dots + \beta_k x_{ik} + \varepsilon_i .
\end{equation}
Empilhando as $n$ observações, $\vetor{y} = \vetor{X}\symbf{\beta} + \symbf{\varepsilon}$, em que $\vetor{X}$ é a matriz $n \times (k+1)$ com uma coluna de uns para o intercepto. O método dos \termo{mínimos quadrados ordinários} (MQO) escolhe os coeficientes que minimizam a soma dos quadrados dos resíduos.

\begin{resultado}{Estimador de mínimos quadrados}
O vetor que minimiza $S(\symbf{\beta}) = (\vetor{y} - \vetor{X}\symbf{\beta})\T(\vetor{y} - \vetor{X}\symbf{\beta})$ é
\begin{equation}
  \hat{\symbf{\beta}} = (\vetor{X}\T\vetor{X})^{-1}\vetor{X}\T\vetor{y} ,
  \label{eq:c09-mqo}
\end{equation}
desde que as colunas de $\vetor{X}$ sejam linearmente independentes. Com uma única feature, a inclinação é $\hat\beta_1 = \widehat{\Cov}(x, y) / \widehat{\Var}(x)$.
\end{resultado}

A Equação~\eqref{eq:c09-mqo} sai de derivar $S$ em relação a $\symbf{\beta}$ e igualar a zero, o que dá $\vetor{X}\T(\vetor{y} - \vetor{X}\hat{\symbf{\beta}}) = \vetor{0}$: os resíduos são ortogonais a todas as features. Cada coeficiente $\hat\beta_j$ é a variação esperada de $y$ quando $x_j$ aumenta uma unidade e as demais features ficam fixas. A qualidade do ajuste é resumida pelo \termo{coeficiente de determinação},
\begin{equation}
  R^2 = 1 - \frac{\sum_i (y_i - \hat y_i)^2}{\sum_i (y_i - \bar y)^2} ,
  \label{eq:c09-r2}
\end{equation}
a fração da variância de $y$ que o modelo reproduz na amostra em que foi estimado.

\subsection{Um exemplo que funciona: o beta de uma ação}

O \termo{modelo de mercado} explica o retorno de uma ação pelo retorno do índice no mesmo dia, $R_{i,t} = \alpha_i + \beta_i R_{m,t} + \varepsilon_{i,t}$ \citep{campbell1997}. O coeficiente $\beta_i$, o beta\index{beta}, mede a sensibilidade do papel ao mercado e já apareceu nos cenários de estresse do Capítulo~\ref{cap:portfolio}. O Código~\ref{cod:c09-beta} o estima para PETR4 contra o Ibovespa com a biblioteca statsmodels, que imprime a tabela clássica de regressão.

\lstinputlisting[style=py, caption={Beta de PETR4 pelo modelo de mercado (\arq{cap09/l01\_beta.py})}, label=cod:c09-beta]{scripts/cap09/l01_beta.py}
\needspace{15\baselineskip}
\begin{lstlisting}[style=saida, basicstyle=\ttfamily\fontsize{7.7}{10.6}\selectfont\color{black!78}]
                            OLS Regression Results
==============================================================================
Dep. Variable:               PETR4.SA   R-squared:                       0.544
Model:                            OLS   Adj. R-squared:                  0.544
No. Observations:                1492   F-statistic:                     1779.
Covariance Type:            nonrobust   Prob (F-statistic):          1.58e-256
==============================================================================
==============================================================================
                 coef    std err          t      P>|t|      [0.025      0.975]
------------------------------------------------------------------------------
const          0.0010      0.000      2.239      0.025       0.000       0.002
^BVSP          1.2671      0.030     42.182      0.000       1.208       1.326
==============================================================================
erro-padrão do beta: MQO 0.0300 | HC0 0.0648 | HAC 0.0749
\end{lstlisting}

A coluna \cod{coef} traz as estimativas. O beta de 1,27 diz que, entre 2019 e 2024, um dia em que o Ibovespa subiu 1\% foi em média um dia em que PETR4 subiu 1,27\%. A coluna \cod{std err} é o erro-padrão de cada estimativa, \cod{t} é a razão entre as duas e \cod{P>|t|} é o valor-p do teste de que o coeficiente verdadeiro é zero. As duas últimas colunas formam o intervalo de 95\% de confiança. O $R^2$ de 0,544 indica que o índice reproduz pouco mais da metade da variância diária da ação; na regressão com uma feature ele é o quadrado da correlação entre as duas séries. A Figura~\ref{fig:c09-beta} mostra a nuvem de pontos e a reta ajustada.

\figura[b]{cap09/beta}{Modelo de mercado para PETR4. À esquerda, cada ponto é um pregão de 2019 a 2024 e a reta é o ajuste de MQO. À direita, o mesmo beta estimado em janelas móveis de 252 pregões: ele passou de 2,0 em 2016 e ficou abaixo de 0,7 em 2025. O coeficiente da regressão é uma média de um parâmetro que não ficou parado.}{fig:c09-beta}

\subsection{As hipóteses e o que falha com dados financeiros}

O estimador da Equação~\eqref{eq:c09-mqo} pode ser calculado para quaisquer dados. As boas propriedades que os livros lhe atribuem, e os erros-padrão que o programa imprime, dependem de hipóteses sobre o erro $\varepsilon$: média zero dadas as features, variância constante (homocedasticidade), ausência de correlação entre erros de datas diferentes e features que não sejam combinação linear umas das outras. \citet{james2021}, no capítulo 3, discutem cada uma. Com retornos, três delas falham de forma previsível.

\begin{itemize}
  \item \textbf{Heterocedasticidade.}\index{heterocedasticidade} A volatilidade se agrupa no tempo (Capítulo~\ref{cap:series}), de modo que a variância do erro em março de 2020 não é a de um mês calmo. Os coeficientes continuam bem estimados em média, mas os erros-padrão usuais ficam errados, quase sempre pequenos demais.
  \item \textbf{Autocorrelação dos erros.} Aparece sempre que o alvo cobre vários dias e as observações são diárias, pela sobreposição da Figura~\ref{fig:c09-linha}. O efeito é o mesmo: uma precisão aparente que não existe.
  \item \textbf{Multicolinearidade.}\index{multicolinearidade} Indicadores técnicos são transformações do mesmo preço e andam juntos. Quando duas features são quase iguais, a matriz $\vetor{X}\T\vetor{X}$ fica próxima de singular e os coeficientes individuais se tornam instáveis, com sinais opostos que se compensam, embora a previsão conjunta pouco mude.
\end{itemize}

Para os dois primeiros problemas existe um remédio padrão: manter os coeficientes de MQO e trocar a fórmula dos erros-padrão por uma que admita variância variável e autocorrelação, a de \citet{newey1987}, conhecida pela sigla HAC. A última linha da saída do Código~\ref{cod:c09-beta} mostra a diferença: o erro-padrão do beta passa de 0,030 para 0,075, duas vezes e meia maior. Quase todo o aumento vem da heterocedasticidade, pois a correção só para ela (HC0) já leva a 0,065; a parcela da autocorrelação, aqui com cinco defasagens, é pequena, como se espera de retornos de um dia. O intervalo de 95\% se alarga de $[1,21;\ 1,33]$ para cerca de $[1,12;\ 1,41]$. A conclusão de que o beta é maior que 1 sobrevive, mas com bem menos folga do que a tabela sugeria. Para a multicolinearidade, as saídas são descartar features redundantes ou penalizar o tamanho dos coeficientes, a regularização do Capítulo~\ref{cap:ml2}.

Há um quarto problema, que nenhuma correção de erro-padrão resolve: o parâmetro muda. O painel direito da Figura~\ref{fig:c09-beta} mostra o beta de PETR4 variando entre 0,64 e 2,03 conforme a janela de um ano escolhida.

\begin{armadilha}{Regressão contemporânea não é previsão}
O $R^2$ de 54\% do modelo de mercado não significa que se consegue prever PETR4. O retorno do Ibovespa no dia $t$ só é conhecido no fechamento de $t$, junto com o da ação. A regressão descreve como os dois se movem juntos, o que serve para medir risco e montar proteções. Para prever, o lado direito precisa estar datado antes do lado esquerdo.
\end{armadilha}

\subsection{A regressão preditiva e o \texorpdfstring{$R^2$}{R2} fora da amostra}

Mude agora uma única coisa: as features passam a ser dez indicadores calculados com dados até o dia $t$ (a Seção~\ref{sec:c09-features} mostra como) e o alvo passa a ser o log-retorno de $t$ a $t+1$. O modelo é estimado com os pregões de 2010 a 2018 e avaliado nos de 2019 a 2025, que ele nunca viu. Dentro da amostra de estimação o $R^2$ da Equação~\eqref{eq:c09-r2} nunca é negativo e nunca diminui quando se acrescentam features, mesmo que sejam ruído. A medida honesta usa os dados de teste e compara o modelo com uma previsão ingênua.

\begin{definicao}{$R^2$ fora da amostra}
Sejam $\hat y_t$ as previsões do modelo para as observações de teste e $\bar y_{\text{ref}, t}$ uma previsão de referência que não usa as features. O \termo{$R^2$ fora da amostra} é
\begin{equation}
  R^2_{\text{fora}} = 1 - \frac{\sum_{t \in \text{teste}} (y_t - \hat y_t)^2}{\sum_{t \in \text{teste}} (y_t - \bar y_{\text{ref}, t})^2} .
  \label{eq:c09-r2oos}
\end{equation}
Ele é positivo se o modelo erra menos que a referência e negativo se erra mais.
\end{definicao}

Duas referências são usuais. \citet{campbell2008} usam a média histórica dos retornos disponível até a data da previsão. \citet{gu2020} usam zero, com o argumento de que a média histórica de uma ação é estimada com tanto ruído que prever zero é uma referência mais difícil de bater. O Código~\ref{cod:c09-preditiva} mostra o núcleo da conta, com a média histórica atualizada a cada dia; a função \cod{base}, apresentada na Seção~\ref{sec:c09-features}, devolve as features e o alvo já alinhados, e \cod{treino} e \cod{teste} marcam as linhas de cada período.

\lstinputlisting[style=py, caption={Regressão preditiva do retorno do dia seguinte (\arq{cap09/l02\_preditiva.py})}, label=cod:c09-preditiva, firstline=9, lastline=20]{scripts/cap09/l02_preditiva.py}
\begin{lstlisting}[style=saida]
treino: 2129 dias | teste: 1741 dias
R2 dentro da amostra:       0.5482%
R2 fora, ref. média hist.:  0.1122%
R2 fora, ref. zero:         0.1133%
R2 fora por ano (ref. zero):
  2019: -0.0429%
  2020:  0.7883%
  2021: -0.1457%
  2022: -0.1222%
  2023:  0.1287%
  2024: -2.9243%
  2025: -0.7310%
\end{lstlisting}

Dentro da amostra, dez indicadores explicam 0,55\% da variância do retorno do dia seguinte. Parece pouco e é menos ainda: quando $k$ features sem qualquer relação com o alvo entram em uma regressão com $n$ observações e erros normais, o $R^2$ esperado é $k/(n-1)$, aqui $10/2128 = 0,47\%$. O ajuste obtido é o que o puro acaso produziria. Fora da amostra o $R^2$ fica em torno de 0,1\%, positivo no agregado, mas negativo em cinco dos sete anos. O número agregado é dominado por 2020, o ano de maior variância do período, que por isso pesa mais nas duas somas da Equação~\eqref{eq:c09-r2oos}. A leitura correta é que este modelo linear, com estas features, não prevê o retorno diário de PETR4 de maneira distinguível de zero.\nota{O método \cod{score} do scikit-learn, aplicado aos dados de teste, usa como referência a média do próprio teste, que não era conhecida na hora da previsão. Prefira escrever a Equação~\eqref{eq:c09-r2oos} explicitamente.}

Resultados dessa ordem são a regra em previsão de retornos. \citet{campbell2008} mostram que valores de $R^2$ fora da amostra muito pequenos já podem ter importância econômica para um investidor, e os números de \citet{gu2020} citados na abertura são da mesma ordem de grandeza. Quem encontra 5\% de $R^2$ fora da amostra em retornos diários deve procurar o vazamento antes de comemorar.

\section{Classificação e regressão logística}

Para decidir entre comprar e ficar de fora, basta acertar o lado do retorno, sem acertar o tamanho. Trocar o alvo contínuo pelo seu sinal transforma a regressão em classificação, com duas vantagens: o alvo deixa de ser dominado por poucos dias extremos, e a saída do modelo passa a ser uma probabilidade, que se converte naturalmente em decisão. A desvantagem é jogar fora a magnitude.

Defina $y_t = 1$ se o retorno do período seguinte for positivo e $y_t = 0$ caso contrário. Ajustar uma reta a um alvo que só vale 0 ou 1 produziria ``probabilidades'' abaixo de zero e acima de um. A \termo{regressão logística} resolve isso passando a combinação linear das features por uma função que comprime qualquer número real para o intervalo entre 0 e 1.

\begin{definicao}{Regressão logística}
O modelo supõe que a probabilidade da classe 1 é
\begin{equation}
  p(\vetor{x}) = \Pr(y = 1 \mid \vetor{x}) = \sigma(\beta_0 + \vetor{x}\T\symbf{\beta}) ,
  \qquad \sigma(z) = \frac{1}{1 + e^{-z}} ,
  \label{eq:c09-logit}
\end{equation}
em que $\sigma$ é a \termo{função sigmoide}. De forma equivalente, o logaritmo da razão de chances é linear nas features:
\begin{equation}
  \ln \frac{p(\vetor{x})}{1 - p(\vetor{x})} = \beta_0 + \vetor{x}\T\symbf{\beta} .
  \label{eq:c09-logodds}
\end{equation}
\end{definicao}

\figura[t]{cap09/logistica}{À esquerda, a função sigmoide, que converte o escore linear $z$ em probabilidade. À direita, a curva ROC da regressão logística do Código~\ref{cod:c09-logistica} nos dados de teste: ela quase coincide com a diagonal, que é o desempenho de um sorteio.}{fig:c09-logistica}

A razão $p/(1-p)$ são as \termo{chances} (\ingles{odds}): uma probabilidade de 60\% corresponde a chances de $0,6/0,4 = 1,5$ para 1. A Equação~\eqref{eq:c09-logodds} dá a interpretação dos coeficientes. Aumentar $x_j$ em uma unidade soma $\beta_j$ ao logaritmo das chances, isto é, multiplica as chances por $e^{\beta_j}$, a \termo{razão de chances}. Um coeficiente de 0,25 multiplica as chances por $e^{0,25} \approx 1,28$. O efeito sobre a probabilidade não é constante: como mostra o painel esquerdo da Figura~\ref{fig:c09-logistica}, ele é máximo perto de $p = 0,5$ e some nas extremidades.

Os coeficientes são estimados por máxima verossimilhança. Para dados $(\vetor{x}_i, y_i)$ independentes, maximizar a verossimilhança equivale a minimizar a \termo{perda logarítmica} (\ingles{log-loss}),
\begin{equation}
  L(\symbf{\beta}) = -\frac{1}{n} \sum_{i=1}^{n} \Big[ y_i \ln p_i + (1 - y_i) \ln (1 - p_i) \Big] , \qquad p_i = p(\vetor{x}_i) ,
  \label{eq:c09-logloss}
\end{equation}
que castiga com força a previsão confiante e errada: atribuir $p_i = 0,99$ a um dia que acabou em queda custa $-\ln 0,01 = 4,6$, contra $-\ln 0,5 = 0,69$ de uma previsão neutra. Não há fórmula fechada como a de MQO; o mínimo é encontrado por um algoritmo iterativo \citep[seção 4.4]{hastie2009}.

A probabilidade vira decisão por meio de um \termo{limiar}: prevê-se a classe 1 quando $p(\vetor{x})$ passa de um valor escolhido, em geral 0,5. O limiar é uma escolha de quem usa o modelo. Subi-lo para 0,55 produz menos sinais de compra, mais seletivos.

\subsection{Medir um classificador}

O Código~\ref{cod:c09-logistica} ajusta a logística às mesmas dez features, com a mesma divisão de treino e teste, e prevê se o dia seguinte será de alta.

\lstinputlisting[style=py, caption={Regressão logística para o sinal do retorno de amanhã (\arq{cap09/l03\_logistica.py})}, label=cod:c09-logistica, firstline=11, lastline=20]{scripts/cap09/l03_logistica.py}
\begin{lstlisting}[style=saida]
matriz de confusão (linhas: real 0, 1):
[[482 348]
 [496 415]]
acurácia  0.5152
precisão  0.5439
recall    0.4555
F1        0.4958
AUC       0.5271
sempre alta (classe majoritária): 0.5233
repetir o sinal de hoje: 0.4957
variação diária média, em módulo: 0.0164
                  ret_1  ret_5  ret_21  vol_21  p_mm50
beta             -0.134  0.080   0.040   0.026   0.135
razão de chances  0.875  1.083   1.041   1.026   1.144
                  rsi_14  macd_h  pct_b  vol_rel  sexta
beta              -0.333  -0.028  0.246    0.069  0.008
razão de chances   0.717   0.972  1.279    1.072  1.008
\end{lstlisting}

A \termo{matriz de confusão} cruza o que aconteceu (linhas) com o que foi previsto (colunas). Dos 1741 dias de teste, o modelo previu alta em $348 + 415 = 763$. Acertou 415, os verdadeiros positivos (VP), e errou 348, os falsos positivos (FP). Dos dias em que previu queda, 482 foram verdadeiros negativos (VN) e 496 foram altas que ele perdeu, os falsos negativos (FN). Todas as métricas usuais saem desses quatro números.

\begin{itemize}
  \item \textbf{Acurácia}, $(\text{VP} + \text{VN})/n = 897/1741 = 51,5\%$: a fração de acertos.
  \item \textbf{Precisão}, $\text{VP}/(\text{VP} + \text{FP}) = 415/763 = 54,4\%$: quando o modelo diz ``alta'', quantas vezes acerta. É a métrica que importa para quem só opera quando há sinal.
  \item \textbf{Recall} (ou sensibilidade), $\text{VP}/(\text{VP} + \text{FN}) = 415/911 = 45,6\%$: de todas as altas que ocorreram, quantas o modelo apontou.
  \item \textbf{F1}, a média harmônica de precisão e recall, $2 \times 0,544 \times 0,456 / (0,544 + 0,456) = 0,496$, que só é alta quando as duas são.
  \item \textbf{AUC}\index{AUC}, a área sob a curva ROC\index{curva ROC}. A curva percorre todos os limiares possíveis e marca, para cada um, a taxa de verdadeiros positivos contra a de falsos positivos. A área tem uma leitura direta: é a probabilidade de o modelo dar nota maior a um dia de alta sorteado ao acaso do que a um dia de queda sorteado ao acaso. Vale 0,5 para um sorteio e 1 para um classificador perfeito, e não depende do limiar escolhido.
\end{itemize}

\begin{armadilha}{Acurácia sem referência não diz nada}
Um modelo que prevê ``alta'' todos os dias, sem olhar feature alguma, acerta 52,3\% do teste, porque essa é a proporção de dias de alta no período. A logística, com 51,5\%, perde para ele. Em problemas com classes desbalanceadas a distorção é pior: se o evento de interesse ocorre em 5\% dos dias, o classificador que nunca o prevê tem 95\% de acurácia e nenhuma utilidade. Sempre compare com a classe majoritária e olhe precisão, recall e AUC.
\end{armadilha}

Os coeficientes impressos referem-se a features padronizadas, de modo que cada razão de chances corresponde ao aumento de um desvio-padrão na feature.\nota{A classe \cod{LogisticRegression} aplica por padrão uma penalização do tipo ridge, com \cod{C=1}. Com 2 mil observações e dez features o efeito é desprezível. O Capítulo~\ref{cap:ml2} trata de regularização.} Os dois maiores em módulo têm sinais opostos: $-0,33$ para o RSI e $+0,25$ para o \%B de Bollinger, duas features cuja correlação na amostra é de 0,89. É o sintoma de multicolinearidade descrito na Seção~\ref{sec:c09-linear}, e seria um erro ler ali uma descoberta sobre o RSI.

\subsection{Quanto vale acertar 53\% das vezes}

Suponha um modelo hipotético que acerte o lado em 53\% dos dias. Essa acurácia parece irrelevante e pode ser a base de um negócio, ou não valer nada. Três contas separam os dois casos.

A primeira é o ganho esperado por aposta. Se acertos e erros tiverem a mesma magnitude média $m$, quem acerta uma fração $p$ das vezes ganha em média $(2p - 1)\,m$ por operação. Com $p = 0,53$ e $m = 1,64\%$, a variação diária média de PETR4 no período de teste, isso dá cerca de 0,1\% por dia, antes de custos. A segunda conta são os custos: a 0,10\% por ponta\index{custo por ponta}, o custo padrão deste livro, o ganho inteiro é consumido se a posição mudar todos os dias. A terceira é a incerteza da própria acurácia. Com $n$ previsões independentes e $p$ perto de 0,5, o erro-padrão da acurácia medida é $\sqrt{0,25/n}$; para os 1741 dias do teste, 1,2 ponto percentual. Mesmo que o modelo hipotético medisse 52,4\% nesse teste, estaria a apenas dois erros-padrão de 50\%, no limite do que o acaso explica sozinho.

A hipótese de magnitudes iguais é a que mais falha: um modelo pode acertar a maioria dos dias pequenos, errar os grandes e perder dinheiro com acurácia acima de 50\%. E um acerto pequeno só vira resultado quando é repetido em muitas apostas pouco correlacionadas, ideia que \citet{grinold2000} formalizam ao relacionar o desempenho de um gestor à qualidade de cada previsão multiplicada pela raiz do número de apostas independentes.

\section{Engenharia de features com dados de mercado}
\label{sec:c09-features}

Chama-se \termo{engenharia de features} o trabalho de transformar dados brutos em variáveis que carreguem informação em um formato que o modelo consiga usar. Com preço e volume diários, as famílias mais comuns estão na Tabela~\ref{tab:c09-features}, e o Código~\ref{cod:c09-features} as implementa para PETR4.

\lstinputlisting[style=py, caption={Construção das features (\arq{cap09/features.py}, trecho)}, label=cod:c09-features, firstline=11, lastline=50]{scripts/cap09/features.py}
\begin{lstlisting}[style=saida]
pregões no arquivo: 3920 | intervalos de mais de 5 dias corridos: 5
linhas completas: 3870, de 2010-03-17 a 2025-12-29
linha de 2020-03-09 igual sem os dados futuros: True
rsi_14  pct_b     0.89
ret_21  p_mm50    0.88
p_mm50  rsi_14    0.85
\end{lstlisting}

\begin{table}[t]
\caption{Features do capítulo, todas calculadas com dados até o fechamento de $t$}
\label{tab:c09-features}
\begin{tabularx}{\textwidth}{>{\ttfamily\small}l >{\raggedright\arraybackslash}X}
\toprule
\cab{Coluna} & \cab{Definição e ideia} \\
\midrule
ret\_1, ret\_5, ret\_21 & Log-retorno acumulado de 1, 5 e 21 pregões. Medem \ingles{momentum} (ou reversão) em vários horizontes. \\
vol\_21 & Desvio-padrão dos log-retornos diários em 21 pregões: a volatilidade realizada recente. \\
p\_mm50 & Preço dividido pela média móvel de 50 dias, menos 1. Distância em relação à tendência. \\
rsi\_14 & Índice de força relativa de Wilder: peso das altas no movimento recente, de 0 a 100. \\
macd\_h & Histograma do MACD: diferença entre médias exponenciais de 12 e 26 dias, menos a sua própria média de 9 dias, em fração do preço. \\
pct\_b & \%B de Bollinger: posição do preço entre as bandas de dois desvios em torno da média de 20 dias (0 na inferior, 1 na superior). \\
vol\_rel & Logaritmo do volume do dia sobre o volume médio de 20 dias. \\
sexta & Indicadora de que $t$ é sexta-feira. Efeito de calendário: o alvo de um dia é então, em geral, o retorno da segunda. \\
\bottomrule
\end{tabularx}
\end{table}

A primeira linha da saída lembra uma consequência da limpeza do Capítulo~\ref{cap:dados}: em cinco pontos de 2011 e 2012 o arquivo salta mais de cinco dias corridos, porque os pregões sem cotação confiável foram removidos. O retorno seguinte a cada lacuna cobre vários pregões e é tratado aqui como se fosse diário. Todos esses pontos ficam no período de treino.

\subsection{RSI, MACD e bandas de Bollinger}

O \termo{índice de força relativa} (RSI, de \ingles{relative strength index}) foi proposto por \citet{wilder1978} e circula em mais de uma versão. Seja $\Delta P_t = P_t - P_{t-1}$. Separe as variações em altas, $U_t = \max(\Delta P_t, 0)$, e quedas, $D_t = \max(-\Delta P_t, 0)$. Wilder suaviza cada série com uma média própria: o primeiro valor é a média simples das $n$ primeiras observações e, a partir daí,
\begin{equation}
  \bar U_t = \frac{(n-1)\,\bar U_{t-1} + U_t}{n} , \qquad
  \bar D_t = \frac{(n-1)\,\bar D_{t-1} + D_t}{n} ,
  \label{eq:c09-wilder}
\end{equation}
com $n = 14$ na versão original. O indicador é
\begin{equation}
  \text{RSI}_t = 100 - \frac{100}{1 + \bar U_t / \bar D_t} = 100\,\frac{\bar U_t}{\bar U_t + \bar D_t} .
  \label{eq:c09-rsi}
\end{equation}
A segunda forma deixa claro o que ele mede: a fração do movimento recente que foi de alta. Valores acima de 70 são lidos, por convenção, como mercado ``sobrecomprado'' e abaixo de 30 como ``sobrevendido''. A suavização da Equação~\eqref{eq:c09-wilder} é uma média exponencial com fator $\alpha = 1/n$, e não $2/(n+1)$, que é a convenção das médias exponenciais comuns. Implementações que usam média móvel simples ou o fator errado produzem outro indicador com o mesmo nome. A função \cod{wilder} do Código~\ref{cod:c09-features} segue a definição original, inclusive na semente.

O \termo{MACD} (\ingles{moving average convergence divergence}) é a diferença entre duas médias móveis exponenciais do preço, uma rápida de 12 dias e uma lenta de 26. Sua própria média exponencial de 9 dias é chamada de linha de sinal, e a diferença entre as duas é o histograma: positivo quando a tendência curta acelera em relação à longa. As bandas de Bollinger, que o Capítulo~\ref{cap:estrategias} usou como regra de reversão, ficam aqui dois desvios-padrão móveis acima e abaixo da média de 20 dias do preço (lá, do logaritmo do preço), e o \%B localiza o preço dentro delas. A Figura~\ref{fig:c09-features} mostra os três em 2024.

\figura[t]{cap09/features}{Três indicadores calculados sobre PETR4 em 2024. No alto, o preço, a média de 20 dias e a faixa das bandas de Bollinger; no meio, o RSI com as linhas de 30 e 70; embaixo, o histograma do MACD. Os três são funções do mesmo preço e contam, em boa parte, a mesma história.}{fig:c09-features}

Os indicadores técnicos são formas diferentes de resumir a trajetória recente do preço, e a saída do Código~\ref{cod:c09-features} confirma a redundância: RSI e \%B têm correlação de 0,89. O que eles oferecem a um modelo são transformações não lineares e já normalizadas dos retornos passados. Se há informação nelas é uma questão empírica.

\subsection{A regra de ouro: só o passado}

Toda feature da linha $t$ deve poder ser calculada no fechamento de $t$, com os dados que existiam naquele momento. No Código~\ref{cod:c09-features} isso é garantido porque \cod{rolling} e \cod{ewm} só olham para trás, e o único deslocamento para a frente, \cod{shift(-h)}, aparece na construção do alvo, na função \cod{base} do mesmo arquivo. A terceira linha da saída é um teste que vale a pena repetir em todo projeto: recalcular as features com a base truncada em uma data e conferir que a linha daquela data não muda.\nota{O preço ajustado muda retroativamente a cada provento (Capítulo~\ref{cap:retornos}). Como um ajuste futuro multiplica todo o passado pelo mesmo fator, features que são razões de preços não se alteram. Uma feature em nível de preço se alteraria.}

\begin{armadilha}{Features que enxergam o futuro}
Os vazamentos mais comuns entram por funções de aparência inocente: \cod{rolling(..., center=True)}, que centraliza a janela e usa dias seguintes; \cod{shift} com sinal trocado; preenchimento de lacunas com \cod{bfill}, que copia o valor futuro para trás; máximos e mínimos calculados sobre a série inteira; dados de balanço datados pelo trimestre a que se referem, e não pelo dia em que foram divulgados. Todos são casos de \ingles{look-ahead}\index{look-ahead} (Capítulo~\ref{cap:backtest}), e um único deles basta para inflar o resultado de um backtest.
\end{armadilha}

\subsection{Estacionariedade e escala}

Um modelo só generaliza se as features do futuro se parecerem com as do treino. O preço em nível não cumpre isso: PETR4 ajustada foi de pouco mais de R\$~1 a mais de R\$~30 na amostra, e um modelo treinado com preços de um dígito nunca viu os valores que encontrará depois. Por isso todas as colunas da Tabela~\ref{tab:c09-features} são retornos, razões ou osciladores limitados, que flutuam em torno de um nível estável.

A transformação tem um custo. O retorno diário é estacionário, mas quase não guarda memória do caminho que o preço percorreu. \citet{lopezdeprado2018}, no capítulo 5, propõe um meio-termo, a \termo{diferenciação fracionária}: em vez de diferenciar a série uma vez inteira, aplica-se uma diferença de ordem $d$ entre 0 e 1, a menor que torne a série estacionária, preservando o máximo de memória. Fica aqui o registro da ideia; não vamos implementá-la.

Features em escalas muito diferentes, como um RSI que vai de 0 a 100 e uma volatilidade diária de 0,02, atrapalham os métodos de otimização e tornam os coeficientes incomparáveis. O remédio é \termo{padronizar}: subtrair a média e dividir pelo desvio-padrão. A média e o desvio precisam ser calculados apenas nos dados de treino e depois aplicados ao teste. Calculá-los na amostra inteira informa ao modelo, ainda que de leve, qual será a volatilidade do futuro. No scikit-learn \citep{pedregosa2011}, o objeto \cod{Pipeline} encadeia o padronizador e o modelo e cuida disso sozinho: cada vez que \cod{fit} é chamado, inclusive dentro de cada dobra de uma validação cruzada, o padronizador é reajustado só com os dados daquele treino. Os classificadores deste capítulo são sempre ajustados dentro de \cod{make\_pipeline} por essa razão; a regressão por MQO do Código~\ref{cod:c09-preditiva} dispensa o cuidado, porque suas previsões não mudam com a escala das features.

\section{Rotulagem: o que exatamente se quer prever}

Em finanças o rótulo é uma decisão de projeto, e muda o problema. O mais simples é o de \termo{horizonte fixo}: o sinal do retorno dos próximos $h$ dias, como na Equação~\eqref{eq:c09-alvo}. \citet{lopezdeprado2018}, no capítulo 3, critica esse rótulo. Dois defeitos interessam aqui. O primeiro é ignorar a volatilidade: um movimento de 2\% em dez dias é grande em um período calmo e irrelevante em uma crise, mas recebe o mesmo rótulo. O segundo é ignorar o caminho. Uma posição que caiu 15\% no terceiro dia e terminou o décimo em leve alta ganha rótulo positivo, embora nenhum operador com limite de perda a tivesse carregado até lá.

O \termo{método da tripla barreira}, do mesmo autor, rotula cada observação pelo que teria acontecido com uma operação de verdade. A partir da entrada em $t$ traçam-se três barreiras: uma superior, de realização de lucro; uma inferior, de limite de perda (\ingles{stop}); e uma vertical, o prazo máximo de $h$ dias. As duas horizontais ficam a uma distância proporcional à volatilidade estimada em $t$. O rótulo é dado pela primeira barreira tocada: $+1$ para a superior, $-1$ para a inferior. Se o prazo vencer antes, o autor discute duas opções, rotular com zero ou com o sinal do retorno acumulado até ali; usamos zero. A Figura~\ref{fig:c09-barreiras} mostra os três desfechos.

\figura[t]{cap09/barreiras}{Tripla barreira em três entradas reais em PETR4, com prazo de 10 pregões e barreiras a um desvio-padrão do retorno de 10 dias. A linha roxa é o caminho até o primeiro toque; a pontilhada, o que veio depois. No painel do meio a perda limite foi atingida no sétimo dia e o preço se recuperou em seguida: pelo horizonte fixo esse rótulo seria positivo.}{fig:c09-barreiras}

\needspace{10\baselineskip}
\lstinputlisting[style=py, caption={Tripla barreira com fechamentos diários (\arq{cap09/l04\_barreiras.py}, trecho)}, label=cod:c09-barreiras, firstline=6, lastline=24]{scripts/cap09/l04_barreiras.py}
\begin{lstlisting}[style=saida]
        fração  dias até o toque
rotulo                          
-1.0     0.228             5.573
 0.0     0.525            10.000
 1.0     0.247             5.861
barreira tocada e retorno de 10 dias com sinal oposto: 102 de 1848
sem rótulo no fim da amostra: 2025-12-15 a 2025-12-30
\end{lstlisting}

No Código~\ref{cod:c09-barreiras}, \cod{sigma} é a volatilidade diária dos 21 pregões até $t$, e a largura das barreiras é $k\sqrt{h}\,\sigma_t$ em log-retorno, a regra da raiz do tempo do Capítulo~\ref{cap:retornos}. Com $k = 1$ e $h = 10$, cerca de metade das entradas termina por prazo, 24,7\% tocam primeiro o alvo e 22,8\% o \ingles{stop}, em média depois de pouco menos de seis pregões. Em 102 das 1848 entradas que tocaram uma barreira horizontal, o retorno de dez dias terminou com o sinal oposto ao do rótulo: são os casos em que o horizonte fixo e a tripla barreira contam histórias contrárias.

Esta implementação simplifica o método em dois pontos: as barreiras são simétricas, quando nada impede distâncias diferentes para lucro e perda, e o toque é verificado só no fechamento, quando máximas e mínimas intradiárias atingiriam as barreiras mais cedo. Note também que as últimas $h$ datas ficam sem rótulo, porque o futuro delas ainda não existe.

\section{Validação temporal}
\label{sec:c09-validacao}

Para estimar o desempenho de um modelo em dados novos, a receita dos livros de aprendizado de máquina é a \termo{validação cruzada} em $k$ dobras (\ingles{k-fold}): embaralham-se as observações, divide-se a amostra em $k$ partes, e cada parte serve uma vez de teste para um modelo treinado nas outras $k-1$ \citep[capítulo 5]{james2021}. O procedimento supõe que as observações são independentes e vêm da mesma distribuição. Com séries financeiras as duas hipóteses falham, e o embaralhamento produz dois vazamentos.

O primeiro é treinar com o futuro. Uma dobra de teste sorteada ao acaso contém dias de 2012 avaliados por um modelo que viu 2020. O segundo é mais sutil e mais grave, e vem da sobreposição de rótulos. Com alvo de 20 dias, o rótulo do dia $t$ e o do dia $t+1$ compartilham 19 dos 20 retornos diários, e as features dos dois dias também são quase idênticas. Se o embaralhamento manda $t$ para o teste e $t+1$ para o treino, o modelo é avaliado em uma pergunta cuja resposta ele praticamente já recebeu.

\begin{intuicao}
Embaralhar uma série com rótulos sobrepostos é aplicar uma prova em que cada questão é uma cópia, com uma palavra trocada, de um exercício resolvido na apostila.
\end{intuicao}

\subsection{Dobras que respeitam o tempo}

A alternativa é nunca treinar com dados posteriores ao teste. A classe \cod{TimeSeriesSplit}\index{TimeSeriesSplit} do scikit-learn faz isso com uma \termo{janela expansível}, a janela ancorada do Capítulo~\ref{cap:estrategias}: a amostra é cortada em $k+1$ blocos consecutivos, e a dobra $i$ treina com os $i$ primeiros blocos e testa no seguinte. A Figura~\ref{fig:c09-dobras} compara os esquemas para 12 observações. Dois parâmetros modificam o desenho básico. Com \cod{max\_train\_size}, o treino fica limitado às observações mais recentes, uma janela deslizante (a rolante do Capítulo~\ref{cap:estrategias}) que esquece o passado distante e faz sentido quando se acredita que o regime muda. Com \cod{gap}, um número fixo de observações entre o treino e o teste é descartado.

\begin{figure}[t]
\centering
\begin{tikzpicture}[x=3.7mm, y=4.4mm, font=\sffamily\footnotesize]
  \newcommand{\linha}[5]{%
    \foreach \i in {#3} \fill[lilasmedio] (#1+\i,#2) rectangle ++(1,0.72);
    \foreach \i in {#4} \fill[roxo] (#1+\i,#2) rectangle ++(1,0.72);
    \foreach \i in {#5} \fill[cinzaclaro] (#1+\i,#2) rectangle ++(1,0.72);
    \foreach \i in {0,...,11} \draw[white, line width=0.8pt] (#1+\i,#2) rectangle ++(1,0.72);
    \draw[line width=0.3pt] (#1,#2) rectangle ++(12,0.72);}
  % (a) k-fold embaralhado
  \node[anchor=south west, inner sep=0pt] at (0,3.1) {\textbf{k-fold embaralhado}};
  \linha{0}{2}{0,1,2,3,5,7,8,9}{4,6,10,11}{}
  \linha{0}{1}{0,3,4,5,6,9,10,11}{1,2,7,8}{}
  \linha{0}{0}{1,2,4,6,7,8,10,11}{0,3,5,9}{}
  % (b) expansível
  \node[anchor=south west, inner sep=0pt] at (16,3.1) {\textbf{TimeSeriesSplit}};
  \linha{16}{2}{0,1,2}{3,4,5}{}
  \linha{16}{1}{0,...,5}{6,7,8}{}
  \linha{16}{0}{0,...,8}{9,10,11}{}
  % (c) max_train_size
  \node[anchor=south west, inner sep=0pt] at (0,-2.4) {\textbf{max\_train\_size=4}};
  \linha{0}{-3.5}{0,1,2}{3,4,5}{}
  \linha{0}{-4.5}{2,3,4,5}{6,7,8}{}
  \linha{0}{-5.5}{5,6,7,8}{9,10,11}{}
  % (d) gap
  \node[anchor=south west, inner sep=0pt] at (16,-2.4) {\textbf{gap=2}};
  \linha{16}{-3.5}{0}{3,4,5}{1,2}
  \linha{16}{-4.5}{0,...,3}{6,7,8}{4,5}
  \linha{16}{-5.5}{0,...,6}{9,10,11}{7,8}
  \foreach \x in {0,16} {
    \draw[-{Stealth[length=1.6mm]}] (\x,-6.1) -- (\x+12,-6.1);
    \node[anchor=north east, inner sep=1pt] at (\x+12,-6.25) {tempo};
    \foreach \y/\n in {2/1, 1/2, 0/3, -3.5/1, -4.5/2, -5.5/3}
      \node[anchor=east, inner sep=2pt, text=cinza] at (\x,\y+0.36) {\n};}
  % legenda
  \fill[lilasmedio] (3.2,-8.6) rectangle ++(1,0.72); \node[anchor=west] at (4.4,-8.24) {treino};
  \fill[roxo] (9.2,-8.6) rectangle ++(1,0.72); \node[anchor=west] at (10.4,-8.24) {teste};
  \fill[cinzaclaro] (15.2,-8.6) rectangle ++(1,0.72); \node[anchor=west] at (16.4,-8.24) {descartado (gap)};
  \draw[line width=0.3pt] (23.6,-8.6) rectangle ++(1,0.72); \node[anchor=west] at (24.8,-8.24) {não usado};
\end{tikzpicture}
\caption{Três dobras sobre 12 observações em quatro esquemas (índices gerados por \arq{cap09/l05\_dobras.py}). No k-fold embaralhado as observações de teste ficam cercadas por vizinhas de treino, anteriores e posteriores. No \cod{TimeSeriesSplit} o treino sempre termina antes do teste; \cod{max\_train\_size} limita o treino ao passado recente e \cod{gap} abre um intervalo entre os dois.}
\label{fig:c09-dobras}
\end{figure}

O \cod{gap} é a ferramenta contra a sobreposição de rótulos. Se o rótulo de $t$ usa retornos até $t+h$, as últimas observações do treino têm rótulos que invadem o período de teste e, o que é o mesmo defeito visto de outro ângulo, que ainda não seriam conhecidos na data em que o teste começa. Para que o último rótulo do treino termine antes de começar o primeiro do teste é preciso descartar $h - 1$ observações. É a purga (\ingles{purging}) apresentada na Seção~\ref{sec:c08-cv}. Em esquemas nos quais o treino também pode vir depois do teste, como o k-fold em blocos, purga-se dos dois lados e ainda se acrescenta o embargo, uma margem extra depois do teste, porque a dependência serial das features faz a contaminação durar um pouco mais que $h$ \citep[capítulo 7]{lopezdeprado2018}.

\subsection{O experimento: o mesmo modelo, quatro notas}

O Código~\ref{cod:c09-validacao} põe os esquemas à prova em um problema com sobreposição forte: prever, a cada dia, se o retorno de PETR4 nos 20 pregões seguintes será positivo (no código, \cod{H = 20}). Dois modelos recebem as mesmas features. A regressão logística tem onze parâmetros e pouca capacidade de memorizar. O segundo modelo, de $k$ vizinhos mais próximos, é o oposto: para classificar um dia, procura os 15 dias de treino com features mais parecidas e responde com a proporção de altas entre eles. Ele entra aqui só como exemplo de modelo flexível, da mesma família dos que o Capítulo~\ref{cap:ml2} estuda.

\lstinputlisting[style=py, caption={Quatro esquemas de validação para o mesmo problema (\arq{cap09/l06\_validacao.py}, trecho)}, label=cod:c09-validacao, firstline=15, lastline=30]{scripts/cap09/l06_validacao.py}
\needspace{13\baselineskip}
\begin{lstlisting}[style=saida]
3851 observações, 54.5% de altas
logística    k-fold embaralhado     AUC 0.542  (dobras: 0.534 a 0.548)
logística    k-fold em blocos       AUC 0.521  (dobras: 0.438 a 0.615)
logística    TimeSeriesSplit        AUC 0.513  (dobras: 0.391 a 0.627)
logística    TimeSeriesSplit + gap  AUC 0.523  (dobras: 0.414 a 0.651)
15 vizinhos  k-fold embaralhado     AUC 0.618  (dobras: 0.576 a 0.661)
15 vizinhos  k-fold em blocos       AUC 0.515  (dobras: 0.438 a 0.571)
15 vizinhos  TimeSeriesSplit        AUC 0.507  (dobras: 0.470 a 0.555)
15 vizinhos  TimeSeriesSplit + gap  AUC 0.504  (dobras: 0.471 a 0.557)
vizinhos a menos de 20 pregões: 1.35 de 15 (ao acaso: 0.15)
vizinhos a 1 pregão: 0.37 de 15
distância mediana: 867 pregões
\end{lstlisting}

\figura[t]{cap09/validacao}{AUC em cada uma das cinco dobras de teste (pontos) e sua média (traço), para os quatro esquemas do Código~\ref{cod:c09-validacao}. Com o k-fold embaralhado, em ocre, as dobras concordam entre si e apontam um modelo com poder de previsão. Nos demais esquemas a média volta para perto de 0,5 e a dispersão entre dobras revela quão incerta é a estimativa.}{fig:c09-validacao}

O modelo de vizinhos recebe AUC de 0,62 no k-fold embaralhado e de 0,50 a 0,51 quando o treino termina antes do teste, com ou sem purga. Nada mudou no modelo nem nos dados; mudou apenas quem estava no treino quando cada dia foi avaliado. As três últimas linhas da saída vêm de um diagnóstico no fim do script, que refaz a busca de vizinhos de cada dobra embaralhada, e mostram o mecanismo. A maior parte dos 15 dias de treino ``mais parecidos'' com um dia $t$ de teste está longe dele no calendário, a uma distância mediana de 867 pregões. Mas em média 1,35 deles cai a menos de 20 pregões de $t$ (0,37 é vizinho imediato), nove vezes mais do que o 0,15 que o acaso colocaria ali, e esses têm rótulos que compartilham retornos com o de $t$. Pouco mais de um vizinho com a resposta quase pronta, em quinze, leva a AUC de 0,50 a 0,62.

A logística sofre menos. Sua média cai de 0,542 para algo entre 0,51 e 0,52, uma diferença pequena diante da dispersão entre dobras, porque um modelo linear com onze parâmetros não tem como decorar vizinhanças. O vazamento por embaralhamento cresce com a flexibilidade do modelo, e por isso é mais perigoso justamente com os modelos do próximo capítulo.

A Figura~\ref{fig:c09-validacao} mostra um segundo engano do embaralhamento. As cinco dobras embaralhadas dão resultados quase iguais, de 0,534 a 0,548 para a logística, o que passa uma impressão de estabilidade. As dobras cronológicas vão de 0,39 a 0,65. Essa dispersão é a incerteza verdadeira: cada bloco de teste tem cerca de 640 dias, que com rótulos de 20 dias equivalem a pouco mais de 30 janelas sem sobreposição. A amostra inteira, de 3851 linhas, contém menos de 200 retornos de 20 dias independentes.

O k-fold em blocos, sem embaralhar, já elimina quase toda a ilusão, porque nele a sobreposição só ocorre nas bordas de cada bloco. Ele continua treinando com o futuro, mas a comparação isola a causa: o problema estava nos vizinhos.

\subsection{Vazamentos típicos e referências obrigatórias}

O embaralhamento é um caso de \termo{vazamento de informação} (\ingles{data leakage}): qualquer caminho pelo qual dados do período de teste influenciam o modelo que será avaliado nele. A Tabela~\ref{tab:c09-vazamentos} reúne os mais frequentes. Todos têm o mesmo efeito, uma métrica otimista que desaparece em produção, e a mesma prevenção: tudo o que é ajustado aos dados, do padronizador à lista de features, fica dentro do laço de validação.

\begin{table}[t]
\caption{Vazamentos típicos em aprendizado de máquina financeiro}
\label{tab:c09-vazamentos}
\begin{tabularx}{\textwidth}{>{\raggedright\arraybackslash}p{30mm} >{\raggedright\arraybackslash}X >{\raggedright\arraybackslash}X}
\toprule
\cab{Vazamento} & \cab{Como acontece} & \cab{Como evitar} \\
\midrule
Feature com futuro & Janela centralizada, \cod{shift} invertido, \cod{bfill}, dado datado antes da divulgação & Teste de truncamento; datar cada dado pelo dia em que ficou disponível \\
Normalização global & Média e desvio calculados na amostra inteira antes de separar o teste & \cod{Pipeline}: o ajuste ocorre só no treino de cada dobra \\
Seleção de features na amostra inteira & Escolher as variáveis mais correlacionadas com o alvo usando todos os dados e só depois validar & Fazer a seleção dentro de cada dobra \citep[seção 7.10]{hastie2009} \\
Hiperparâmetros ajustados no teste & Experimentar configurações até a métrica de teste melhorar & Três conjuntos: treino, validação e um teste final usado uma vez (Capítulo~\ref{cap:ml2}) \\
Rótulos sobrepostos & Alvo de vários dias com treino e teste vizinhos & Purga de $h-1$ observações e embargo (Seção~\ref{sec:c08-cv}) \\
\bottomrule
\end{tabularx}
\end{table}

O quarto item é o que menos parece trapaça. Cada vez que você olha o resultado no teste e volta para mudar o modelo, o teste passa a fazer parte do treino, como no sobreajuste de backtest dos Capítulos~\ref{cap:backtest} e~\ref{cap:estrategias}.

Por fim, todo modelo deve ser comparado, no mesmo teste, com pelo menos duas referências: a classe majoritária (ou a previsão zero, em regressão) e uma regra simples, de uma linha. No Código~\ref{cod:c09-logistica}, ``sempre alta'' acertou 52,3\% e ``repetir o sinal de hoje'' acertou 49,6\%. Um modelo que não supera as duas não aprendeu nada que valha a sua complexidade.

\section{Da previsão à estratégia}

Entre uma probabilidade e uma ordem há uma regra de posição, que diz quanto comprar para cada valor do sinal, e um backtest com custos, nos moldes dos Capítulos~\ref{cap:backtest} e~\ref{cap:estrategias}. O Código~\ref{cod:c09-backtest} percorre esse caminho da forma mais simples possível. No início de cada ano, de 2019 a 2025, a logística é reajustada com todos os dados anteriores e gera probabilidades para o ano que começa, um esquema \ingles{walk-forward}\index{walk-forward} de janela expansível. A posição $w_t$ vale 1 (comprado em PETR4) se a probabilidade de alta passa de 0,5 e 0 (aplicado no CDI) caso contrário. O retorno da estratégia entre os fechamentos de $t$ e $t+1$ é
\begin{equation}
  R_{t+1} = w_t\,R^{\text{PETR4}}_{t+1} + (1 - w_t)\,R^f_{t+1} - c\,\lvert w_t - w_{t-1} \rvert ,
  \label{eq:c09-estrategia}
\end{equation}
com $R^f$ a taxa diária do CDI\index{CDI} (arquivo \arq{dados/c09\_cdi.csv}, gravado por \arq{cap09/baixar\_cdi.py}) e $c = 0,10\%$ o custo por ponta, a hipótese única do livro, que resume taxas, metade do \ingles{spread} e \ingles{slippage} de ordens pequenas; não inclui o impacto de mercado de ordens grandes, tratado no Capítulo~\ref{cap:trading}. A convenção é a de \emph{execução no fechamento do sinal}\index{execucao no fechamento do sinal@execução no fechamento do sinal}, a mesma do motor do Capítulo~\ref{cap:backtest}: a posição é assumida no próprio fechamento de $t$ que gerou o sinal. Não há \ingles{look-ahead}, porque o retorno do dia $t$ não é capturado, mas a hipótese é otimista e o código a testa repetindo a conta com um pregão de atraso. A referência é \termo{comprar e manter} (\ingles{buy and hold}): adquirir a ação no primeiro dia e não negociar mais.

\figura[t]{cap09/backtest}{Resultado fora da amostra da estratégia guiada pela regressão logística, de 2019 a 2025, contra comprar e manter PETR4 e contra o CDI. Sem custos o modelo fica perto da ação; com 0,10\% por ponta fica bem abaixo dela; com um pregão de atraso na execução perde para o CDI.}{fig:c09-backtest}

\lstinputlisting[style=py, caption={Do sinal à posição e ao resultado líquido (\arq{cap09/l07\_backtest.py}, trecho)}, label=cod:c09-backtest, firstline=19, lastline=38]{scripts/cap09/l07_backtest.py}
\needspace{8\baselineskip}
\begin{lstlisting}[style=saida]
modelo, sem custo     23.2% a.a.  Sharpe  0.55  DD máx. -45.2%
modelo, com custo     12.6% a.a.  Sharpe  0.26  DD máx. -45.6%
com custo e atraso     2.8% a.a.  Sharpe -0.05  DD máx. -63.5%
comprar e manter      24.8% a.a.  Sharpe  0.54  DD máx. -63.4%
CDI no período: 9.1% a.a.
dias comprado: 48.9% | pontas: 620 | acerto do sinal: 50.7%
\end{lstlisting}

As taxas são anualizadas pela convenção de 252 dias (Seção~\ref{sec:anualizar}), o índice de Sharpe usa o retorno em excesso ao CDI, e a última coluna é o drawdown máximo do Capítulo~\ref{cap:risco}. Antes de custos, a estratégia rendeu 23,2\% ao ano, um pouco menos que os 24,8\% de comprar e manter, com Sharpe praticamente igual (0,55 e 0,54). As duas verificações seguintes tiram o que restava. As 620 pontas negociadas custam, somadas, 62 pontos percentuais de retorno em sete anos, e o retorno líquido cai para 12,6\% ao ano, com Sharpe de 0,26. Com um pregão de atraso na execução, o retorno cai para 2,8\% ao ano, abaixo do CDI de 9,1\%. O sinal acerta 50,7\% dos dias e troca de lado a cada três pregões: não há previsão explorável aqui.\nota{O Capítulo~\ref{cap:backtest} obtém Sharpe de 0,56 para comprar e manter PETR4 em 2019--2025. A diferença vem da amostra: lá entra o retorno do pregão de 2 de janeiro de 2019; aqui a primeira posição é assumida no fechamento desse dia.}

O drawdown máximo menor, de 45\% contra 63\%, tem uma explicação que dispensa o modelo: a carteira passou metade dos dias fora da ação. Na versão com atraso ele é igual ao da própria ação.

Duas ressalvas valem para qualquer leitura desses números. O ativo foi escolhido hoje, sabendo que a Petrobras atravessou o período e pagou dividendos elevados; é o viés de sobrevivência da Seção~\ref{sec:c07-sobrevivencia}, e 24,8\% ao ano não é o que se espera de uma ação tomada ao acaso em 2019. E o teste cobre um único papel em um único intervalo.

O resultado é decepcionante e é o esperado: indicadores técnicos diários sobre uma das ações mais negociadas do país são informação que todos os participantes têm. O que fica do exercício é o encanamento, que agora existe e não vaza: features datadas, padronização dentro do \ingles{pipeline}, validação cronológica com purga, referências ingênuas, custos e teste de atraso. É nele que o Capítulo~\ref{cap:ml2} vai ligar modelos mais flexíveis.

\begin{resumo}
  \item Em mercados a razão sinal-ruído é mínima, as relações mudam e as observações dependem umas das outras. Um $R^2$ fora da amostra de décimos de ponto percentual é o resultado normal, e métricas muito melhores pedem primeiro uma busca por vazamento.
  \item A regressão linear por MQO funciona bem para descrever relações contemporâneas, como o beta. Com retornos, os erros-padrão usuais são otimistas e devem ser trocados por versões HAC; o próprio coeficiente muda com o tempo.
  \item A regressão logística transforma features em probabilidade por meio da sigmoide. Avalie com matriz de confusão, precisão, recall e AUC, sempre ao lado da classe majoritária e de uma regra simples.
  \item Toda feature da data $t$ usa apenas dados até $t$, em forma estacionária, e qualquer transformação ajustada aos dados, como a padronização, é estimada só no treino.
  \item O rótulo é uma escolha: o horizonte fixo ignora volatilidade e caminho, e a tripla barreira os incorpora. Rótulos de vários dias se sobrepõem e reduzem a informação efetiva da amostra.
  \item Embaralhar séries temporais vaza o futuro para o treino. Use divisões cronológicas, purgue $h-1$ observações entre treino e teste e só acredite em uma previsão depois de vê-la sobreviver a um backtest com custos.
\end{resumo}

\needspace{7\baselineskip}
\begin{exercicios}
  \item \nivel{básico} Um classificador produz, em 200 dias de teste, 60 verdadeiros positivos, 50 falsos positivos, 50 falsos negativos e 40 verdadeiros negativos. Calcule acurácia, precisão, recall e F1, e a acurácia de quem prevê sempre a classe majoritária.
  \item \nivel{básico} Uma regressão logística tem $\beta_0 = 0,1$ e um único coeficiente $\beta_1 = -0,4$ para uma feature padronizada. Calcule a probabilidade prevista para $x = 0$ e para $x = 2$, e a razão de chances associada a um desvio-padrão de aumento em $x$.
  \item \nivel{intermediário} Escreva uma versão do RSI que use média móvel simples no lugar da suavização de Wilder e outra que use \cod{ewm(span=14)}. Compare as três séries para PETR4 em 2024: qual é a maior diferença, em pontos, e em quantos dias elas discordam sobre o indicador estar acima de 70?
  \item \nivel{intermediário} Introduza de propósito um vazamento no Código~\ref{cod:c09-features}, trocando a média móvel de 50 dias por uma versão com \cod{center=True}. Refaça o Código~\ref{cod:c09-logistica} e relate o que acontece com a acurácia e a AUC. O teste de truncamento detecta o erro?
  \item \nivel{intermediário} No Código~\ref{cod:c09-validacao}, varie o horizonte do rótulo para 1, 5, 10 e 40 dias e registre a diferença de AUC entre o k-fold embaralhado e o \cod{TimeSeriesSplit} com \cod{gap}, para os dois modelos. Explique o padrão que aparece.
  \item \nivel{desafio} Use os rótulos da tripla barreira do Código~\ref{cod:c09-barreiras}, descartando os de prazo, como alvo de uma regressão logística. Valide com \cod{TimeSeriesSplit}, escolhendo o \cod{gap} adequado, e compare com a classe majoritária em cada dobra.
  \item \nivel{desafio} Refaça o backtest do Código~\ref{cod:c09-backtest} para os sete ativos brasileiros de \arq{precos\_ajustados.csv}, usando apenas as features que não dependem de volume. Em quantos deles a estratégia líquida supera comprar e manter em índice de Sharpe? O que esse número diria se fosse quatro de sete?
\end{exercicios}

\begin{leituras}
  \item \citet{james2021}, capítulos 3, 4 e 5, para regressão linear, classificação e validação cruzada em nível introdutório.
  \item \citet{hastie2009}, capítulos 3, 4 e 7, com o tratamento formal dos mesmos temas; a seção 7.10 discute o jeito errado e o jeito certo de fazer validação cruzada.
  \item \citet{lopezdeprado2018}, capítulos 3, 5 e 7, para tripla barreira, diferenciação fracionária e validação com purga e embargo; \citet{lopezdeprado2020}, seção 5, retoma a rotulagem em texto mais curto.
  \item \citet{jansen2020}, capítulos 4, 6 e 7, com muitos exemplos em Python de construção de features e de modelos lineares aplicados a retornos.
\end{leituras}
